\documentclass[aspectratio=169,10pt]{beamer}
\usetheme[numbering=fraction,progressbar=frametitle]{metropolis}
\usepackage{fontspec}
\setsansfont{Latin Modern Sans}
\usepackage{amsmath,amssymb,mathtools,graphicx,booktabs}
\usepackage{polyglossia}
\setdefaultlanguage{english}
\setbeamertemplate{navigation symbols}{}
\setbeamerfont{frametitle}{size=\large}
\setbeamerfont{normal text}{size=\normalsize}
\setlength{\parskip}{.45em}
\hypersetup{pdfauthor={},pdfsubject={Lecture slides},pdfkeywords={},pdftitle={Graphs}}
\providecommand{\N}{\mathbb{N}}
\providecommand{\Z}{\mathbb{Z}}
\providecommand{\D}{\mathbb{D}}
\providecommand{\Q}{\mathbb{Q}}
\providecommand{\R}{\mathbb{R}}
\providecommand{\C}{\mathbb{C}}
\providecommand{\vect}[1]{\overrightarrow{#1}}
\providecommand{\vi}{\vec{\imath}}
\providecommand{\vj}{\vec{\jmath}}
\providecommand{\vk}{\vec{k}}
\providecommand{\norm}[1]{\left\lVert #1\right\rVert}
\providecommand{\vc}[2]{\begin{pmatrix}#1\\#2\end{pmatrix}}
\providecommand{\deter}[2]{\det\left(#1,#2\right)}
\providecommand{\abs}[1]{\left\lvert #1\right\rvert}
\providecommand{\dd}{\,\mathrm{d}}
\providecommand{\e}{\mathrm{e}}
\providecommand{\Card}{\operatorname{Card}}
\providecommand{\E}{\operatorname{E}}
\providecommand{\V}{\operatorname{V}}
\providecommand{\degre}{^{\circ}}
\providecommand{\un}[1]{\,\mathrm{#1}}
\providecommand{\sci}[2]{#1\times10^{#2}}
\providecommand{\pgcd}{\operatorname{PGCD}}
\providecommand{\ppcm}{\operatorname{PPCM}}
\title{Graphs}
\subtitle{Lecture and worked examples}
\date{}
\author{}
\begin{document}
\begin{frame}[plain]
\vfill
{\usebeamerfont{title}\inserttitle\par}
\medskip
{\usebeamerfont{subtitle}\insertsubtitle\par}
\medskip
{\usebeamercolor[fg]{progress bar}\rule{\textwidth}{.5pt}}
\vfill
\end{frame}
\begin{frame}{Objectives}
Model a situation with a graph; use the vocabulary of graphs (order, degree, walk, cycle, connectedness, complete graph, directed graph); write the adjacency matrix of a graph; count the walks of given length with the powers of this matrix.
\end{frame}
\begin{frame}{Introductory activity}
Five towns \(A\), \(B\), \(C\), \(D\), \(E\) are linked by direct bush-taxi routes, in both directions: \(A\)-\(B\), \(A\)-\(C\), \(A\)-\(D\), \(B\)-\(C\), \(C\)-\(D\) and \(D\)-\(E\).
\end{frame}
\begin{frame}{Introductory activity (2)}
\begin{enumerate}
\setcounter{enumi}{0}\item Represent the network by a diagram: one point per town, one line per route.

\item For each town, count the number of routes leaving it. Compute the sum of these numbers and compare it with the number of routes. Explain.

\item List the journeys from \(A\) to \(C\) that use exactly two routes (a town may be visited again). How many are there?

\item Same question for the journeys from \(A\) to \(E\) using exactly three routes.

\end{enumerate}
\end{frame}
\begin{frame}{Introductory activity (3)}
\begin{enumerate}
\setcounter{enumi}{4}\item We store the information in a table \(M\) with \(5\) rows and \(5\) columns: the entry in row \(X\) and column \(Y\) is \(1\) if there is a direct route between \(X\) and \(Y\), \(0\) otherwise. Write \(M\), then compute \(M^2\) and \(M^3\) with a calculator. Compare some entries with questions 3 and 4.

\item State a conjecture.

\end{enumerate}
\end{frame}
\begin{frame}{Vocabulary}
A (undirected) \emph{graph} consists of a finite set of \emph{vertices} and a set of \emph{edges}, each edge joining two distinct vertices; two joined vertices are \emph{adjacent}. We consider graphs without multiple edges. The \emph{order} of the graph is its number of vertices. The \emph{degree} of a vertex is the number of edges of which it is an end. A graph is \emph{complete} if any two distinct vertices are adjacent.
\end{frame}
\begin{frame}{Vocabulary}
\begin{columns}[T,onlytextwidth]
\begin{column}{.57\textwidth}
A \emph{walk} is a sequence of vertices in which any two consecutive vertices are adjacent; its \emph{length} is its number of edges (a vertex or an edge may be used again). A walk is \emph{closed} if its ends are equal. A \emph{cycle} is a closed walk of length at least \(3\) whose edges are distinct. A graph is \emph{connected} if any two vertices are joined by a walk.
\end{column}
\begin{column}{.4\textwidth}
\centering
\includegraphics[width=\linewidth,height=4.5cm,keepaspectratio]{figures/expertes-10-1-en.pdf}
\end{column}
\end{columns}
\end{frame}
\begin{frame}{Vocabulary}
In a graph, the sum of the degrees of the vertices is twice the number of edges. Consequently, the number of vertices of odd degree is even.
\end{frame}
\begin{frame}{Vocabulary}
Each edge has two ends: it is counted twice in the sum of the degrees. The sum of the degrees is even; now a sum of integers is even if and only if it contains an even number of odd terms.
\end{frame}
\begin{frame}{Vocabulary}
The complete graph of order \(n\) has \(\binom n2=\dfrac{n(n-1)}2\) edges.
\end{frame}
\begin{frame}{Vocabulary}
In a \emph{directed} graph, the links, called \emph{arcs}, have a direction: the arc \((X,Y)\) goes from \(X\) to \(Y\). A \emph{path} is a sequence of vertices in which each is joined to the next by an arc. A graph is \emph{weighted} when each edge or arc carries a number, called its weight.
\end{frame}
\begin{frame}{Adjacency matrix}
The vertices of a graph of order \(n\) being numbered from \(1\) to \(n\), its \emph{adjacency matrix} is the square matrix \(M\) of order \(n\) whose entry \(m_{ij}\) is \(1\) if the vertices \(i\) and \(j\) are joined (by an edge, or by an arc from \(i\) to \(j\) for a directed graph), and \(0\) otherwise. For an undirected graph, \(M\) is symmetric: \(m_{ij}=m_{ji}\).
\end{frame}
\begin{frame}{Adjacency matrix}
Let \(M\) be the adjacency matrix of a graph and \(n\ge1\) an integer. The entry in row \(i\) and column \(j\) of \(M^n\) equals the number of walks (of paths, for a directed graph) of length \(n\) from vertex \(i\) to vertex \(j\).
\end{frame}
\begin{frame}{Adjacency matrix}
Let \(m^{(n)}_{ij}\) be the entries of \(M^n\) and \(c^{(n)}_{ij}\) the number of walks of length \(n\) from \(i\) to \(j\). For \(n=1\), a walk of length \(1\) from \(i\) to \(j\) is an edge: \(c^{(1)}_{ij}=m_{ij}\).
\end{frame}
\begin{frame}{Adjacency matrix (2)}
Assume \(m^{(n)}_{ij}=c^{(n)}_{ij}\) for all \(i\), \(j\). A walk of length \(n+1\) from \(i\) to \(j\) consists of a walk of length \(n\) from \(i\) to a vertex \(k\), followed by an edge from \(k\) to \(j\). Sorting these walks by the second-to-last vertex \(k\), we obtain
\[c^{(n+1)}_{ij}=\sum_{k=1}^{p}c^{(n)}_{ik}\,m_{kj}=\sum_{k=1}^{p}m^{(n)}_{ik}\,m_{kj}=m^{(n+1)}_{ij},\]
where \(p\) is the order of the graph, since \(M^{n+1}=M^nM\). The property holds for \(n+1\).
\end{frame}
\begin{frame}{Adjacency matrix}
The diagonal entry \(\left(M^2\right)_{ii}\) is the degree of vertex \(i\) (undirected graph). A graph of order \(p\ge2\) is connected if and only if all the off-diagonal entries of \(M+M^2+\dots+M^{p-1}\) are positive: a walk between two distinct vertices can always be shortened into a walk with no repeated vertex, of length at most \(p-1\).
\end{frame}
\end{document}
